% ---------------------------------------------------------------------------
%  QML Lab research project 01: Quantum-Kernel Intrusion Detection
%  Build:  latexmk -pdf main.tex   (needs qmlproject.sty, harvardqml.bst, references.bib)
% ---------------------------------------------------------------------------
\documentclass[11pt,a4paper]{article}
\usepackage{qmlproject}
\qmlheader{Project 01: Quantum-kernel IDS}

\hypersetup{pdftitle={QML Lab project 01: Quantum-Kernel Intrusion Detection}}

\begin{document}

\begin{center}
  {\small\textcolor{black!55}{QML Lab \textbullet\ research project 01}}\\[0.5em]
  {\LARGE\bfseries Quantum-Kernel Intrusion Detection}\\[0.6em]
  {\large Can a quantum kernel separate normal from malicious network traffic better than a tuned classical kernel?}\\[0.6em]
  {\small Living document \textbullet\ last updated \today}
\end{center}

\begin{projectcard}
\small
\begin{tabularx}{\linewidth}{@{}>{\bfseries}l X >{\bfseries}l X@{}}
Quadrant  & \CQ\ (classical data, quantum processor) & Track & Kernels \\
Difficulty & Medium & Tools & Qiskit, scikit-learn \\
Builds on & Modules 4 and 6 of the QML Lab & Workspace & \path{qml_projects/quantum-kernel-intrusion-detection/} \\
\end{tabularx}
\end{projectcard}

\paragraph{Summary.}
This project applies a quantum-kernel support vector machine to network intrusion detection, a
security task with real, imbalanced data. Traffic records are reduced to a few features with PCA,
encoded into qubits with a ZZ feature map, and classified with an SVM that uses the state overlap
as its kernel. The comparison that matters is against an RBF kernel and a small neural network
with the same features. Kernel bandwidth, the scaling of features before encoding, is treated as
the main control knob, because it decides whether the quantum kernel generalises or simply memorises.

% ===========================================================================
\section{Background and key concepts}
% ===========================================================================
\subsection{Intrusion detection as classification}
A network intrusion detection system labels connections or flows as normal or as one of several
attack types. Public benchmarks make results comparable. NSL-KDD \citep{tavallaee2009} removes the
duplicate records that inflated results on the older KDD Cup 99 data and has 41 features per
connection and four attack families (DoS, Probe, R2L, U2R). CICIoT2023 \citep{neto2023} is a newer
data set recorded on a real IoT network, with many attack types and flow-level features. Both are
strongly imbalanced: some attack classes are rare, so accuracy alone is misleading.

\subsection{Support vector machines and kernels}
An SVM \citep{cortes1995} learns a maximum-margin separator. In its dual form it only needs the
kernel $k(x,x')$ between pairs of training points:
\begin{equation}
  \max_{\alpha}\ \sum_i\alpha_i-\frac12\sum_{i,j}\alpha_i\alpha_jy_iy_j\,k(x_i,x_j)
  \quad\text{s.t.}\quad 0\le\alpha_i\le C,\ \ \sum_i\alpha_iy_i=0,
\end{equation}
and predicts $f(x)=\operatorname{sign}\big(\sum_i\alpha_iy_ik(x_i,x)+b\big)$. Any positive
semi-definite similarity can be plugged in, including one computed by a quantum circuit.

\subsection{The quantum kernel}
The ZZ feature map of \citet{havlicek2019} encodes $x\in\R^d$ on $d$ qubits with Hadamards and
phases that depend on single features and on pairwise products. The kernel is the fidelity of two
encoded states,
\begin{equation}
  \kappa(x,x')=|\braket{\psi(x)|\psi(x')}|^2=|\braket{0|U_\Phi(x')^\dagger U_\Phi(x)|0}|^2,
\end{equation}
estimated on hardware as the probability of measuring all zeros after running
$U_\Phi(x')^\dagger U_\Phi(x)$. Building the Gram matrix needs $O(N^2)$ circuit evaluations for
$N$ training points, which limits training sets to a few thousand points.

\subsection{Bandwidth and concentration}
Before encoding, features are scaled as $x\mapsto\lambda x$. For large $\lambda$ the encoded states
are almost orthogonal, the Gram matrix approaches the identity and the model memorises the training
set; this is exponential concentration \citep{thanasilp2024}. \citet{shaydulin2022} showed that
reducing the bandwidth $\lambda$ restores generalisation, in the same way as the width of an RBF
kernel. \citet{huang2021power} give a test for whether a quantum kernel can beat classical ones at
all on given data: the \emph{geometric difference}
\begin{equation}
  g(K^{C}\Vert K^{Q})=\sqrt{\big\lVert\sqrt{K^{Q}}\,(K^{C})^{-1}\sqrt{K^{Q}}\big\rVert_\infty},
\end{equation}
which must be large for a quantum advantage to be possible.

% ===========================================================================
\section{Research questions}
% ===========================================================================
The project is open-ended: these questions are the starting point, not the finish line. Add new
ones to the README and the journal as they come up.

\begin{rqlist}
  \item At the best bandwidth, does the quantum-kernel SVM beat a tuned RBF-SVM and a small MLP on binary detection with the same 4--8 features?
  \item Does the quantum kernel help on rare attack classes (per-class recall), or only on the easy majority classes?
  \item How do accuracy, kernel concentration and the geometric difference change with bandwidth $\lambda$ and qubit number $d$?
\end{rqlist}

% ===========================================================================
\section{Method}
% ===========================================================================
\subsection{Data}
NSL-KDD first (KDDTrain+ for training, KDDTest+ for testing), then CICIoT2023. Categorical
features are one-hot encoded, all features standardised, and PCA reduces them to $d\in\{4,6,8\}$
components, scaled to $[0,\pi]$ and then multiplied by $\lambda$. Because the Gram matrix costs
$O(N^2)$ circuit runs, training uses stratified subsamples of $N=250$ to $2000$ connections; the
test set keeps its original class balance.

\subsection{Models}
Fidelity kernel with the ZZ feature map (\texttt{reps} 1--2), simulated exactly with state vectors;
the precomputed Gram matrix goes into \texttt{sklearn.svm.SVC} \citep{pedregosa2011} with $C$
chosen by cross-validation. Multi-class detection uses one-vs-rest.

\subsection{Baselines}
RBF-SVM with $(C,\gamma)$ tuned on the same features; logistic regression; an MLP with a similar
number of parameters; and, as a sanity check, the quantum kernel with $\lambda\to0$, which
approaches a constant kernel.

\subsection{Metrics}
\begin{itemize}
  \item Accuracy, macro-F1, attack recall and false-positive rate; ROC-AUC for the binary task.
  \item Per-class recall for the four NSL-KDD attack families.
  \item Kernel diagnostics: off-diagonal variance, kernel--target alignment and geometric difference $g$.
  \item Cost: number of circuit evaluations and wall-clock time.
\end{itemize}

% ===========================================================================
\section{Planned experiments}
% ===========================================================================
Each experiment becomes one \texttt{[[experiments]]} entry in the workspace, so its runs, parameters
and metrics are stored and compared in the portal.

\begin{table}[h]
\centering\small
\begin{tabularx}{\linewidth}{@{}l>{\raggedright}X>{\raggedright}X>{\raggedright\arraybackslash}X@{}}
\toprule
ID & Experiment & What varies & Main metrics \\
\midrule
  E1 & Binary detection, $d=4$ & bandwidth $\lambda\in\{0.05,\dots,1\}$ & macro-F1, off-diagonal variance \\
  E2 & Qubit scaling & $d\in\{4,6,8\}$ at the best $\lambda$ & macro-F1, $g$, run time \\
  E3 & Attack families & one-vs-rest over DoS, Probe, R2L, U2R & per-class recall \\
  E4 & Learning curves & $N\in\{250,500,1000,2000\}$ & test macro-F1 vs RBF \\
  E5 & Shot noise & kernel entries from $M$ shots & accuracy loss vs exact kernel \\
\bottomrule
\end{tabularx}
\caption{Experiment plan. Results go into the journal and the results log (\cref{sec:log}).}
\end{table}

% ===========================================================================
\section{Running it in the QML Lab}
% ===========================================================================
The workspace lives in \path{qml_projects/quantum-kernel-intrusion-detection/}. Scripts import three helpers from
\path{qml_lab/plots.py}:
\begin{itemize}
  \item \texttt{params()} reads the parameters from the portal's run form (defaults when run by hand);
  \item \texttt{save()} stores figures with the run;
  \item \texttt{log\_metrics()} records numbers that appear in the run table.
\end{itemize}

\begin{lstlisting}
qml_projects/quantum-kernel-intrusion-detection/
  project.toml          metadata + experiments
  README.md             goal, open questions, ideas
  experiments/          runnable scripts
  docs/main.tex         this document
\end{lstlisting}

First experiment: This experiment is \emph{planned}: write the script, add the block below to \texttt{project.toml}, then run \texttt{python manage.py sync\_projects}.

\begin{lstlisting}[language={}]
[[experiments]]
slug = "nsl-kdd-binary"
title = "Binary NSL-KDD detection with a ZZ quantum kernel"
script = "experiments/e01_nsl_kdd_binary.py"
figure = "kernel_vs_rbf.png"

[experiments.params]
n_features = 4
bandwidth = 0.5
reps = 2
n_train = 1000
seed = 0
\end{lstlisting}

\begin{lstlisting}[language=bash]
# from the portal folder: pick up changes to project.toml
python manage.py sync_projects
cd qml_projects/quantum-kernel-intrusion-detection
# run by hand with the default parameters
PYTHONPATH=../../qml_lab python experiments/e01_nsl_kdd_binary.py
\end{lstlisting}

% ===========================================================================
\section{Quality bar}
% ===========================================================================
\begin{itemize}
  \item A classical baseline with a similar number of trainable parameters, tuned with the same care \citep{bowles2024}.
  \item At least five random seeds per configuration, reported as mean $\pm$ standard deviation.
  \item How results change with the number of qubits.
  \item At least one run under a realistic noise model.
  \item Report per-class recall and false-positive rate, not accuracy alone; the data is imbalanced.
\end{itemize}

% ===========================================================================
\section{Open directions}
% ===========================================================================
\begin{itemize}
  \item Move from NSL-KDD to CICIoT2023 and compare how conclusions change on recent IoT traffic.
  \item One-class SVM with the quantum kernel for anomaly detection when attack labels are missing.
  \item Compare inference cost and accuracy with a TinyML-sized classical detector running on an edge device.
  \item Projected quantum kernels (Project 03) as a less concentrated alternative.
\end{itemize}

% ===========================================================================
\section{Risks and limitations}
% ===========================================================================
\begin{itemize}
  \item $O(N^2)$ kernel evaluations limit training-set size; subsampling may favour simpler models.
  \item PCA to 4--8 components can discard the signal of rare attacks; check what the baselines lose too.
  \item NSL-KDD has known artefacts and is old; treat it as a benchmark, not as evidence of real-world performance.
  \item Large bandwidth or many qubits leads to concentrated, memorising kernels.
\end{itemize}

% ===========================================================================
\section{Results log}
\label{sec:log}
% ===========================================================================
Summarise findings here as they become stable; day-to-day notes belong in the portal journal.

\begin{table}[h]
\centering\small
\begin{tabularx}{\linewidth}{@{}l l >{\raggedright}X >{\raggedright\arraybackslash}X@{}}
\toprule
Date & Run & Finding & Next step \\
\midrule
 & & & \\
 & & & \\
 & & & \\
\bottomrule
\end{tabularx}
\end{table}

\bibliographystyle{harvardqml}
\bibliography{references}

\end{document}
